% This is the ADASS_template.tex LaTeX file, 24 Oct 2025.
% It is based on the ASP general author template file, but modified to reflect the specific
% requirements of the ADASS proceedings.

% This is the aspauthor.tex LaTeX file
% Copyright 2021, Astronomical Society of the Pacific Conference Series
% Revision:  29 November 2021

% To compile, at the command line positioned at this folder, type:
% pdflatex ADASS_template
% pdflatex ADASS_template
% This will create a file called ADASS_template.pdf

\documentclass[11pt,twoside]{article}

% Do NOT use ANY packages other than asp2023. 
\usepackage{asp2023}

\aspSuppressVolSlug
\resetcounters

% References must all use BibTeX entries in a .bibfile.
% References must be cited in the text using \citet{} or \citep{}.
% Do not use \cite{}.
% See ManuscriptInstructions.pdf for more details
\bibliographystyle{asp2021}

% The ``markboth'' line sets up the running heads for the paper.
% 1 author: "Surname"
% 2 authors: "Surname1 and Surname2"
% 3 authors: "Surname1, Surname2, and Surname3"
% >3 authors: "Surname1 et al."
% Replace ``Short Title'' with the actual paper title, shortened if necessary.
% Use mixed case type for the shortened title
% Ensure shortened title does not cause an overfull hbox LaTeX error
% See ASPmanual2010.pdf 2.1.4  and ManuscriptInstructions.pdf for more details
\markboth{Jeschke, Neugarten and Kackley}{Transitioning Subaru's Remote Observing Software to a Web Site}

\begin{document}

\title{Transitioning Subaru's Remote Observing Software to a Web Site}

% Note the position of the comma between the author name and the 
% affiliation number.
% Authors surnames should come after first names or initials, eg John Smith, or J. Smith.
% Author names should be separated by commas.
% The final author should be preceded by "and".
% Affiliations should not be repeated across multiple \affil commands. If several
% authors share an affiliation this should be in a single \affil which can then
% be referenced for several author names. If only one affiliation, no footnotes are needed.
% See ManuscriptInstructions.pdf and ASP's manual2010.pdf 3.1.4 for more details
\author{Eric~Jeschke,$^1$ Andrew~Neugarten,$^1$ and Russell~Kackley$^1$}
\affil{$^1$Subaru~Telescope,~National~Astronomical~Observatory~of~Japan}

% This section is for ADS Processing.  There must be one line per author. paperauthor has 9 arguments.
\paperauthor{Eric~Jeschke}{eric@naoj.org}{0000-0001-9586-4482}{Subaru~Telescope,~National~Astronomical~Observatory~of~Japan}{Software Division}{Hilo}{Hawaii}{96720}{USA}
\paperauthor{Andrew~Neugarten}{aneugart@naoj.org}{}{Subaru~Telescope,~National~Astronomical~Observatory~of~Japan}{Software Division}{Hilo}{Hawaii}{96720}{USA}
\paperauthor{Russell~Kackley}{rkackley@naoj.org}{}{Subaru~Telescope,~National~Astronomical~Observatory~of~Japan}{Software Division}{Hilo}{Hawaii}{96720}{USA}

% There should be one \aindex line (commented out) for each author. These are used to
% build up the author index for the Proceedings. The surname must come first, followed by
% initials. Note the use of ~ before each initial to control spacing.
% The \author entries (see above) have surname last. These \aindex entries have
% surname first.
% The Aindex.py command willl create them for you after you have constructed the \author
% The first entry should be the first author, for bold-facing the author index page-reference
%\aindex{Jeschke,~E.}
%\aindex{Neugarten,~A.}
%\aindex{Kackley,~R.}


\begin{abstract}
  A number of observatories use VNC (Virtual Network Computing) servers and
  clients to monitor the observation screens from open-use observers
  connecting from remote locations. Usually this involves a custom software
  package that must be installed and configured with some sort of security key
  and/or password to secure the system.  This puts an onus on software
  developers to support multiple platforms (Linux, Mac OS X and Windows)
  for the software, and additionally, provide support for installation issues,
  connection issues and assorted remote troubleshooting.

  We present a case study in transitioning from a remote software package to
  a secure web site with second factor authentication at Subaru Telescope.
  The site is supported by all modern browsers on all platforms, and provides
  excellent viewing capabilities for all screens---without requiring any
  intrusive changes on the observatory software side.  This transition has
  allowed us to minimize installation hassles for remote observers and
  greatly simplify troubleshooting and maintenance because there is no longer
  any custom remote software for the user to install.  We discuss the
  architecture of the new system, what pieces needed to be developed for the
  observatory side and provide some limited anecdotal data on use (the system
  went into use in August 2026).
\end{abstract}

% These lines show examples of subject index entries. At this stage these have to commented
% out, and need to be on separate lines. Eventually, they will be automatically uncommented
% and used to generate entries in the Subject Index at the end of the Proceedings volume.
% Don't leave these in! - replace them with ones relevant to your paper.
%\ssindex{FOOBAR!conference!ADASS 2019}
%\ssindex{FOOBAR!organisations!ASP}

% These lines show examples of ASCL index entries. At this stage these have to commented
% out, and need to be on separate lines. Eventually, they will be automatically uncommented
% and used to generate entries in the ASCL Index at the end of the Proceedings volume.
% The ascl.py command will scan your paper on possible code names.
% Don't leave these in! - replace them with ones relevant to your paper.
%\ooindex{FOOBAR, ascl:1101.010}

\section{Introduction}
Subaru Telescope is an 8.2 meter class optical and infrared telescope situated on
the summit of Maunakea, on the island of Hawai`i.
Until the advent of the 2020 COVID pandemic, Subaru Telescope only provided for
remote observation from two dedicated locations: Hilo, Hawai`i and Mitaka, Japan.
Observers would come to one of those locations to observe during classical
observations, which at the time constituted a significant portion of all observation
time (unattended queue observations have increased since then, but still only account for
a fraction of all observations at Subaru).

\subsection{The Motivation}
With the advent of the pandemic in March 2020, the limitations of this policy became
starkly apparent: observers could join observations using teleconferencing only, without
being able to monitor the observation screens.  The observatory could share one screen,
or point the videocamera at a batch of screens, but neither of these solutions were
very satisfying.

The Software Division began considering how to securely
allow observations from other locations (basically anywhere with ``reasonable''
internet service).
From a technical point of view, this was not conisdered a very difficult problem,
as the observation control system user interfaces are displayed in a series of VNC
(Virtual Network Computing) servers, which are then connected to by VNC clients
running at the summit, Hilo and Mitaka.  VNC security is not particularly
robust, however, and we wanted to control carefully exactly \emph{who} could connect
and \emph{when}.

\section{First iteration: Python + SSH + TigerVNC}
The first iteration of our Gen2 Extended Remote System (GERS) was rolled out quickly
(May 2020) and used a small Python package hosted on Github to establish SSH encrypted
connections to tunnel the VNC connections via a host straddling the Subaru
firewall.
As seen in Figure \ref{ex_fig1},
reverse-forwarding tunnels were established from the observation control system
hosts and standard forwarding tunnels from the remote user's host.
The user would then start a VNC client on their computer connected to a
\emph{local} port and the forwarding ensured that they got connected to the
corresponding VNC server screen at the summit.

\articlefigure[width=0.65\textwidth]{P807_f1}{ex_fig1}{Architecture of the the first iteration of GERS.}

Users would generate and upload a public SSH key and we used a scheduling system
(shown as ``GERS update daemon'') to rotate the keys in and out of the accounts based on the
observation schedule.
In this system the first authentication factor was the key, and the second factor was a TOTP
code.
A GERS administration web site and a user self-service site were established for the
staff and remote users, respectively, to handle the scheduling and uploading/downloading
of keys and configuration files.  A relational database is used to store the credentials,
schedule and public keys.

\subsection{Support Issues}
Supporting a desktop-type package in the current era you need to be able
to handle Windows, Mac OS, and a few flavors of Linux.  We referred users to the
download site for the TigerVNC client and to Conda for a reasonable cross-platform
way to run the Python app.  Many astronomers are now familiar with installing Conda and
running Python applications, but we still found a number of users to have trouble
installing the program, running the program or getting the VNC client to install or
work correctly.  These issues required occasional remote troubleshooting, sometimes
just before their observation.

\section{Second iteration: Python + SSH + web}
After a couple of semesters running with the first iteration, we made a change to the
system to use a web application called noVNC\citep{noVNC}.
noVNC allows you to view a VNC window within a web page using a HTML5 standards-compliant
browser---no VNC client needed.  Figure \ref{ex_fig2} shows the changes.

By adapting the Python program to start up a small local-only web server,
we could now have the users open up a ``menu page'' from which they could click to
open any one of the 8 VNC windows as browser tabs.  Several windows can be viewed
at once, depending on their desktop real-estate, by moving tabs to new browser windows.
This allowed us to drop the need (and remote support) for installing
and running a VNC client; users just needed to have a ``reasonable'' web browser installed,
such as Firefox, Chrome, Safari, Edge, etc.  We still used the Python
program to handle SSH tunneling and the keys to enforce the time periods in which
each observer could connect.  This change reduced the support burden somewhat,
since we no longer had to troubleshoot getting a VNC client installed and running.

\articlefigure[width=0.65\textwidth]{P807_f2}{ex_fig2}{Architecture of the the second iteration of GERS.}

%% The other addition we made during this period was a VNC fanout program (``vnc\_fanout''),
%% which sits between the actual VNC server and the noVNC websocket server---this allows
%% a single VNC connection to be multiplexed between a number of clients,
%% which reduces the burden on the actual VNC servers of being innundated with too many
%% connections.

\section{Third iteration: web only}
The third iteration of the GERS system was rolled out in the fall semester of 2026,
and does away entirely with the need for the user to download and install any software.
We accomplished this by switching from noVNC to \emph{Guacamole}\citep{Guacamole}---a fully
web-based system offered under the auspices of the Apache Software Foundation.
Refer to Figure \ref{ex_fig3} for the following discussion.  

Guacamole offers HTTPS for security and a control database that allowed us to register
users and schedule when they could connect to the system.
This obviated the need to upload SSH keys and swap them in and out of a fixed number
of SSH logins, so the GERS update daemon is changed to simply reading the GERS schedule
from the GERS database and updating the Guacamole database.
In addition to the user/password login on the web site, we ask the users to download and install a
\emph{client} web certificate in their browser, which acts as a second factor authentication.

\articlefigure[width=0.65\textwidth]{P807_f3}{ex_fig3}{Architecture of the third iteration of GERS.}

\subsection{Benefits and Summary}
One of the limitations of the first two iterations was that the user had only a limited
window in which to verify their connection (typically a few hours before the start of
their observation).  This late check often led to last-minute troubleshooting by our
engineers.  Under the Guacamole system we can control which screens the user has access
to by date, so we can allow them to connect up to a week before their observation to
a test screen; this is sufficient to show that they have connected successfully.
On the night of their observation they are able to see any of the 8 observation screens.

Moving to this third iteration has reduced the support required for the GERS remote
observation, as there is no software to download or install and no need to generate and
upload an SSH key.  Security is improved on our firewall interface where we no longer need to
support tunneled forwarded connections, but only secure http.  Furthermore, the Guacamole
VNC viewer running in the browser is full featured, and allows arbitrary scaling
of screens to account for different sizes of monitor that each user may be using.
GERS remote observations are now our default method of observing and the system is very
popular with observers.



%%%%%%%%%%%%%%%%%%%%%%%%%%%%%%%%%%%%%%%%

\section{References}

\bibliography{P807}  % For BibTex

% if we have space left, we might add a conference photograph here. Leave commented for now.
% \bookpartphoto[width=1.0\textwidth]{foobar.eps}{FooBar Photo (Photo: Any Photographer)}

\end{document}
